%==============================================================================
\section{Introduction}
\label{sec:introduction}
%==============================================================================

Sorted collections are fundamental to systems software.
Priority queues, leaderboards, sorted result sets, MIDI note timelines,
and particle simulation indices all require efficient maintenance of order
as underlying data changes.
The standard design choice is between two poles:

\begin{itemize}
  \item \textbf{Flat sorted arrays.} $O(1)$ indexed access, cache-optimal sequential
    scans, SIMD-friendly comparisons~\cite{drepper2007memory}.
    Update cost is $O(n \log n)$ for a full re-sort.
  \item \textbf{Dynamic ordered structures (B-tree, red-black tree).}
    $O(\log n)$ per update. Sequential access costs $O(n)$ traversal;
    rank queries cost $O(r)$ per query or $O(n)$ to materialize.
\end{itemize}

Production systems have converged toward trees for write-heavy workloads
(databases use B$^+$-trees; language standard libraries offer
\texttt{std::map}, \texttt{BTreeMap}), and toward arrays for read-heavy or
batch-processed collections.
When updates are small and \emph{localized}---affecting $k \ll n$ elements
per batch---neither choice is obviously optimal.
The sorted array requires an expensive full re-sort; the B-tree pays $O(\log n)$
per update but makes range iteration or rank queries costly.

\paragraph{Update-aware sorting.}
A recent line of work~\cite{deltasort} observes that in many practical systems,
the \emph{positions} of updated elements are known at the time of sorting
(they were just written by the application).
Given this information, \DeltaSort~\cite{deltasort} restores sorted order of
a flat array in $O(n\sqrt{k})$ expected time and $O(k)$ space,
outperforming both BIS ($O(nk)$) and ESM ($O(n + k\log k)$) for $k \lesssim 1\%$.

\paragraph{Our question.}
The original \DeltaSort paper establishes the algorithm's superiority
\emph{within} the sorted-array family (vs.\ BIS and ESM).
We ask a different question: \emph{does update-aware sorting shift the
trade-off between flat arrays and B-trees, and if so, in which real workloads?}

\paragraph{Findings.} We find:

\begin{enumerate}
  \item \textbf{The crossover is nearly universal.}
    For any workload with even a single indexed read per update epoch
    ($Q \geq 1$), a sorted array maintained by \DeltaSort outperforms
    a \texttt{BTreeMap} for all $(n, k)$ pairs we test.
    The break-even reads-per-epoch $\Qbreak$ is at most $2$ across
    $n \in [10^3, 10^6]$, $k \in [1, 10^4]$, predicted exactly by a
    simple analytic formula (\secref{sec:crossover}).

  \item \textbf{MIDI sequencing is a natural fit.}
    A music sequencer stores events sorted by timestamp, receives small
    batches of edits (drag, quantize, paste), and answers sorted reads
    continuously.
    \DeltaSort over a sorted \texttt{Vec} delivers
    $5$--$20\times$ throughput over \texttt{BTreeSet} for drag and
    quantize operations at $n \leq 10^5$ (\secref{sec:midi}).

  \item \textbf{Large-scale HPC and 3D rendering are \emph{not} good fits.}
    In cosmological particle simulations at $n = 10^6$, \ESM dominates
    \DeltaSort for $k/n \geq 5\%$ because the linear merge pass is
    more cache-efficient than DeltaSort's segment scan.
    In 3D game-engine rendering (Bevy), the phase sort is $<5\%$ of
    the frame budget even under worst-case fully-sorted conditions,
    making algorithmic choices irrelevant to FPS (\secref{sec:negative}).
\end{enumerate}

These results paint a coherent picture: \DeltaSort is most valuable in
\textbf{medium-$n$, low-$k$, read-inclusive} workloads where the collection
is too large for full re-sort to be cheap but too read-heavy for a B-tree's
traversal cost to be acceptable.
MIDI editing, leaderboard maintenance, and interactive dashboards naturally
fall into this regime.
Large-batch HPC and GPU-bound rendering do not.

\paragraph{Artifact.} All benchmarks are implemented in Rust and are
available at \url{https://github.com/shudv/deltasort}.
